\chapter{Images Without a Daemon: Jib and the Local Registry}
\label{ch:jib}

\section{Why not just \code{docker build} in CI?}

The Dockerfiles work. Three reasons the pipeline builds images
differently:

\begin{enumerate}
  \item \textbf{CI needs no Docker daemon.} \code{docker build} requires
  a running daemon; giving Jenkins one means docker-in-docker or mounting
  the host's Docker socket --- which is root on the host, handed to every
  build. Jib builds and pushes images in pure Java: no daemon, no
  privilege.
  \item \textbf{Layering that matches how code changes.} The Dockerfile
  copies one fat jar: one huge layer, fully re-pushed per build. Jib
  splits the application into layers by volatility --- dependencies /
  resources / classes --- so a one-line code change re-pushes only the
  small classes layer. On two cores and Wi-Fi this is the difference
  between seconds and minutes per service.
  \item \textbf{OpenShift-ready by accident.} Jib images run as a
  non-root user with group-0 file permissions --- exactly what OpenShift's
  arbitrary-UID security model demands. Portability for free.
\end{enumerate}

\section{The configuration, annotated}

Jib is configured once in the parent \code{pom.xml}'s
\code{pluginManagement}; each runnable module opts in with a bare
\code{<plugin>} reference, and \code{common} opts \emph{out} with
\code{jib.skip=true}:

\begin{lstlisting}[language=xml]
<plugin>
  <groupId>com.google.cloud.tools</groupId>
  <artifactId>jib-maven-plugin</artifactId>
  <version>3.4.4</version>
  <configuration>
    <from><image>eclipse-temurin:25-jre-alpine</image></from>  (*@\co{1}@*)
    <to><image>${docker.registry}/${project.artifactId}:${project.version}</image></to> (*@\co{2}@*)
    <allowInsecureRegistries>true</allowInsecureRegistries>    (*@\co{3}@*)
    <container>
      <mainClass>${start-class}</mainClass>                    (*@\co{4}@*)
      <creationTime>USE_CURRENT_TIMESTAMP</creationTime>       (*@\co{5}@*)
    </container>
  </configuration>
</plugin>
\end{lstlisting}

\begin{description}
  \item[\co{1}] Same base as the Dockerfiles --- runtime behavior
  unchanged, only the packaging path differs.
  \item[\co{2}] Resolves to \code{localhost:5000/course-service:1.0.0-SNAPSHOT}.
  The name is written \emph{from the cluster's point of view} ---
  Section~\ref{sec:whylocalhost}.
  \item[\co{3}] The local registry is plain HTTP; Jib refuses insecure
  pushes unless told. Every tool in this chain (Jib, containerd) makes
  the same secure-by-default choice and needs the same explicit opt-out.
  \item[\co{4}] Set explicitly because Jib's bundled bytecode scanner
  does not yet parse Java~25 class files --- a real example of
  bleeding-edge Java taxing the toolchain.
  \item[\co{5}] Jib defaults to epoch-0 timestamps for reproducible
  builds; this trades that away so \code{docker images} shows honest
  build times.
\end{description}

\section{The registry, and why \code{localhost:5000}}
\label{sec:whylocalhost}

\begin{lstlisting}[language=bash]
docker run -d --restart=unless-stopped --name registry \
  -p 5000:5000 -v registry-data:/var/lib/registry registry:2
\end{lstlisting}

The non-obvious fact that forces this architecture: \textbf{k3s does not
see Docker's images.} k3s runs its own containerd with its own image
store; \code{docker build} on the host puts images somewhere k3s never
looks. A registry is the honest interface between builder and cluster
--- which the pipeline wants anyway.

An image name is not just a name --- it is \emph{where to pull from}:
\code{localhost:5000/course-service} means ``pull from
localhost:5000.'' Since the puller is the k3s node itself, and pushing
Jenkins shares the host's network namespace (Chapter~22), the same
string is valid for both --- and \code{localhost} can never drift the way
a LAN IP or hostname can. Two configuration hooks make it work:

\begin{lstlisting}[language=yaml]
# /etc/rancher/k3s/registries.yaml -- containerd assumes HTTPS otherwise
mirrors:
  "localhost:5000":
    endpoint: ["http://localhost:5000"]
\end{lstlisting}

\begin{pitfall}[http: server gave HTTP response to HTTPS client]
The signature error of this chapter. Both Jib (pushing) and containerd
(pulling) default to TLS; \code{registry:2} speaks plain HTTP. If a pod
sticks in \code{ImagePullBackOff} with this text in
\code{kubectl describe pod}, the \code{registries.yaml} entry is missing
--- and note containerd reads it \emph{only at startup}:
\code{systemctl restart k3s} after editing.
\end{pitfall}

\begin{decision}[local registry vs.\ GHCR]
GitHub Container Registry would make images pullable anywhere and add
vulnerability scanning --- for the price of pushing hundreds of megabytes
up residential Wi-Fi on every build. A LAN-local registry keeps the
build loop at disk speed. The decision flips the moment a second machine
must pull images; nothing in the manifests would change except the
image prefix.
\end{decision}

\section{Tags, digests, and the SNAPSHOT trap}

A \emph{tag} (\code{:1.0.0-SNAPSHOT}, \code{:a1b2c3d}) is a movable
pointer; a \emph{digest} (\code{@sha256:...}) is the immutable content
hash. The pipeline pushes every build twice-labelled: the SNAPSHOT tag
(which the manifests reference) and the git-SHA tag (which deploys pin
--- Chapter~23). Kubernetes caches images by tag; a re-pushed SNAPSHOT
would be silently ignored on the node --- which is why every Deployment
sets \code{imagePullPolicy: Always}: re-check the registry each start;
if the digest matches, the check costs one manifest fetch.

\begin{quiz}
\begin{enumerate}
  \item Name the three CI problems Jib solves over
  daemon-based builds, and which one is a security issue.
  \item Why must the image name be written from the \emph{cluster's}
  viewpoint, and what breaks if Jenkins and k3s disagree about what
  \code{localhost:5000} means?
  \item A colleague suggests \code{docker save | k3s ctr images import}
  instead of a registry. What does that cost per deploy, and which
  pipeline stage stops working naturally?
  \item Explain how a re-pushed \code{:1.0.0-SNAPSHOT} plus default pull
  policy produces ``I deployed but nothing changed,'' and both fixes.
\end{enumerate}
\end{quiz}
